\documentclass[11pt]{article}
\usepackage[utf8]{inputenc}
\usepackage{amsmath,amssymb}
\usepackage{hyperref}
\title{Scalable Aerosol Cloud Interaction for Large-Scale Settings}
\author{Anonymous}
\date{}
\begin{document}
\maketitle

\begin{abstract}
This paper studies Aerosol Cloud Interaction. Grounded in a real, citable literature \cite{ref1,ref2,ref3,ref4,ref5,ref6,ref7,ref8},
we propose a focused method and report \textbf{illustrative} experiments.
All references below are real and verifiable (OpenAlex/arXiv); the reported
numbers are simulated to illustrate intended behaviour.
\end{abstract}

\section{Introduction}
Aerosol Cloud Interaction is an active research area. This work targets a concrete gap identified from
the recent literature and contributes a method together with an evaluation protocol.

\section{Related Work (real literature)}
The following works, retrieved from public bibliographic APIs, frame our study:
\begin{itemize}
\item Seinfeld et al. (1998) — "<i>Atmospheric Chemistry and Physics: From Air Pollution to Climate Change</i>", Physics Today (cited 9021).
\item Bond et al. (2013) — "Bounding the role of black carbon in the climate system: A scientific assessment", Journal of Geophysical Research Atmospheres (cited 6781).
\item Prior work (1998) — "Atmospheric chemistry and physics: from air pollution to climate change", Choice Reviews Online (cited 5172).
\item Kanakidou et al. (2005) — "Organic aerosol and global climate modelling: a review", Atmospheric chemistry and physics (cited 3739).
\item Strang (1968) — "On the Construction and Comparison of Difference Schemes", SIAM Journal on Numerical Analysis (cited 3509).
\item Andreae et al. (2008) — "Aerosol–cloud–precipitation interactions. Part 1. The nature and sources of cloud-active aerosols", Earth-Science Reviews (cited 1879).
\end{itemize}

\section{Method}
We decompose the problem across shards with a communication-efficient aggregation rule that keeps overhead sublinear in scale. We instantiate this idea for Aerosol Cloud Interaction and analyse its cost/quality trade-off.

\section{Experiments (Illustrative)}
\begin{center}
\begin{tabular}{lcc}
System & Primary Metric & Efficiency \\ \hline
Strong baseline & 0.812 & 1.00$\times$ \\
Ablation & 0.828 & 0.92$\times$ \\
\textbf{Ours} & \textbf{0.851} & \textbf{0.71$\times$}
\end{tabular}
\end{center}
\emph{Metrics simulated to illustrate the intended effect; not measured.}

\section{Conclusion}
We connected a real literature to a concrete method for Aerosol Cloud Interaction. Future work will
validate the illustrative results with real experiments.

\begin{thebibliography}{9}
\bibitem{ref1} John H. Seinfeld, Spyros Ν. Pandis, K. Noone. <i>Atmospheric Chemistry and Physics: From Air Pollution to Climate Change</i>. \emph{Physics Today}, 1998. DOI:10.1063/1.882420
\bibitem{ref2} Tami C. Bond, Sarah J. Doherty, D. W. Fahey et al.. Bounding the role of black carbon in the climate system: A scientific assessment. \emph{Journal of Geophysical Research Atmospheres}, 2013. DOI:10.1002/jgrd.50171
\bibitem{ref3} . Atmospheric chemistry and physics: from air pollution to climate change. \emph{Choice Reviews Online}, 1998. DOI:10.5860/choice.35-5721
\bibitem{ref4} Maria Kanakidou, John H. Seinfeld, Spyros Ν. Pandis et al.. Organic aerosol and global climate modelling: a review. \emph{Atmospheric chemistry and physics}, 2005. DOI:10.5194/acp-5-1053-2005
\bibitem{ref5} Gilbert Strang. On the Construction and Comparison of Difference Schemes. \emph{SIAM Journal on Numerical Analysis}, 1968. DOI:10.1137/0705041
\bibitem{ref6} Meinrat O. Andreae, Daniel Rosenfeld. Aerosol–cloud–precipitation interactions. Part 1. The nature and sources of cloud-active aerosols. \emph{Earth-Science Reviews}, 2008. DOI:10.1016/j.earscirev.2008.03.001
\bibitem{ref7} Yoram J. Kaufman, D. Tanré, L. A. Remer et al.. Operational remote sensing of tropospheric aerosol over land from EOS moderate resolution imaging spectroradiometer. \emph{Journal of Geophysical Research Atmospheres}, 1997. DOI:10.1029/96jd03988
\bibitem{ref8} David S. Lee, D. W. Fahey, Agnieszka Skowron et al.. The contribution of global aviation to anthropogenic climate forcing for 2000 to 2018. \emph{Atmospheric Environment}, 2020. DOI:10.1016/j.atmosenv.2020.117834
\end{thebibliography}
\end{document}
